\documentclass[11pt]{article}
\usepackage[a4paper,margin=1.05in]{geometry}
\usepackage{amsmath,amssymb,booktabs,graphicx,xcolor}
\usepackage[expansion=false,protrusion=true]{microtype}
\usepackage[hidelinks]{hyperref}
\usepackage{caption,natbib,url,enumitem,titlesec,array}
\captionsetup{font=small,labelfont=bf,width=.95\textwidth}
\titleformat{\section}{\normalfont\large\bfseries}{\thesection}{0.6em}{}
\titleformat{\subsection}{\normalfont\normalsize\bfseries}{\thesubsection}{0.6em}{}
\setlist{nosep,leftmargin=1.2em}
\newcommand{\sm}{\textsc{SkillMerge}}
\setlength{\tabcolsep}{5pt}
\renewcommand{\arraystretch}{1.08}

\title{\vspace{-1.6em}\textbf{Task-Scoped Composition of Agent Skills:\\ Three Strategies Measured Against Loading the Whole Skill}\vspace{-0.3em}}
\author{Mathieu Bourgault\thanks{Code, the demo corpus, the results and the scripts that
regenerate every number in this article: \url{https://github.com/mathieubrglt/skillmerge}. The
production corpus is not redistributed; see Section~\ref{sec:avail}.}}
\date{26 August 2026}

\begin{document}
\maketitle

\begin{abstract}
\noindent An agent that needs three skills loads three whole documents, most of which do not apply
to the task in front of it. We built a service that indexes a skill library once and returns, per
request, a single composite scoped to one task and capped at a token budget, and we tested three
ways of representing the library inside it. On a corpus we generated, with redundancy planted by
construction, composition matched whole-skill loading within measurement error at 42\% of the
context ($\Delta = -0.022$ on a blind rubric, 95\% CI $[-0.061, +0.018]$, $p = 0.30$). On 45 real
skills in the Agent Skills format the same method lost 7.5 rubric points and 11 of 12 tasks
($p_{\text{Holm}} = 0.029$), and its cross-skill de-duplication was rejected by an independent
adjudication panel on 56 of 60 sampled merges. Splitting each unit into a general \emph{obligation}
and the specific \emph{lesson} that makes it actionable, and merging only the former while keeping
every one of the latter, raised merge precision from 0.067 to 0.817 (difference $+0.750$, 95\% CI
$[+0.634, +0.866]$, $p < 0.001$) without improving the deliverables. An ablation shows why: printing
the obligations costs quality, because a capable model already holds them. We retain the third
representation and reject its rendering. The case for composition rests on context cost and library
hygiene; no variant matched whole-skill loading on output quality.
\end{abstract}

\section{Introduction}

Skill composition does not improve the work an agent produces. Across three strategies, 12 tasks
and 132 blind-graded deliverables, no composite came within 8 rubric points of loading the skills a
practitioner would have chosen. It does reduce injected context by 46\% to 60\%, and one of the
three strategies makes a skill library de-duplicable without losing content. Those two results, and
the reasons the quality gap persists, are what this article reports.

An agent skill is a document, typically a \texttt{SKILL.md}, that tells a model how a particular
kind of work is done \citep{anthropic2025skills}. Skills are written to stand alone so they can be
authored and shared independently. Two costs follow. A task drawing on three skills pays three
times for whatever they share, and each skill is written for a class of situations, so most of what
is loaded does not apply to the one at hand; models attend less reliably to material buried in long
context \citep{liu2024lost}.

We built \sm{}, a service exposed over the Model Context Protocol \citep{anthropic2024mcp} that
indexes a library once and returns a task-scoped composite. All model-dependent work happens at
index time, so the request path is deterministic and involves no model call. Section~\ref{sec:method}
describes what is common to the three strategies and how they were evaluated.
Sections~\ref{sec:s1} to \ref{sec:s3} report the strategies in the order they were built, because
each was a response to a measured failure of the one before. Section~\ref{sec:retained} states what
we keep.

\section{Related work}

Retrieval-augmented generation \citep{lewis2020rag} established retrieving passages rather than
whole documents. The move here is the same one applied to procedural rather than factual context,
with two differences: the corpus is small, curated and written by hand, and the retrieved units
must be assembled into something an agent can follow as instructions, which makes ordering and
de-duplication part of the problem.

Prompt compression shrinks a prompt that already exists. LLMLingua \citep{jiang2023llmlingua} and
its successors \citep{pan2024llmlingua2} drop low-information tokens with a small language model,
below the semantic level and on the request path. \sm{} selects whole competencies above the
semantic level and off the request path; a composite could then be compressed, so the two compose.

Our offline generation of predicted trigger phrasing is document expansion by query prediction
\citep{nogueira2019doc2query} applied to procedural fragments, and it is what makes a purely
lexical request path workable. Ranking uses BM25 \citep{robertson2009bm25}. Diversity reranking
\citep{carbonell1998mmr} appears in Strategy~1 and is removed in Section~\ref{sec:s1abl} because it
did not earn its place.

\section{Method}
\label{sec:method}

\subsection{Two corpora}

\textbf{Synthetic (Strategy~1).} Measuring de-duplication requires knowing what is duplicated, so we
built a corpus from a ground-truth fragment library: 34 competencies in the earth-observation and
commodity-intelligence domain, of which 10 were planted in five or six skills each, 12 in two to
four, and 12 in one. Shared competencies were written in three paraphrases apiece, and 12 skills
were rendered with per-skill headings, scaffolding prose and shuffled section order, so recovering
the structure is an inference problem rather than an inversion of the generator. The result is 12
skills, 99 fragment instances over 34 distinct competencies, 10,137 tokens.

\textbf{Real (Strategies~2 and 3).} 45 skills exactly as their authors wrote them: 10 from an engineering plugin, 8
document-production skills shipped as platform defaults, and 27 example skills covering consumer
tasks, design and tooling. They total 85,905 tokens, median 999 and maximum 9,921. Twelve ship
auxiliary files, 200 of them. None was written with composition in mind.

\subsection{Tasks, rubrics and grading}

Each study uses a 12-task suite of realistic work requests, with the skills a practitioner would
load recorded as an oracle set. The synthetic study's tasks sit in its own domain. The real study's
suite spans two genres: seven engineering tasks (review a pull request, work a partial-failure
incident, write an architecture decision record, judge release readiness, prioritise technical debt,
chase a silent data-corruption bug, write a runbook) and five document-production tasks (tracked
changes in a long contract, a deck from a report against a template, cleaning a malformed export,
extracting line items from scanned invoices, planning an MCP server).

Every cell of every experiment was produced by an independent agent that saw only its prompt file,
did not know which condition it was in or that other conditions existed, and received identical
length guidance. Grading was blind: all candidates for a task were anonymised, shuffled and scored
together on each rubric criterion at $\{0, 1, 2\}$, plus a 1 to 10 judgement of how useful the
response would be to the person who has to act on it. Two independent panels graded every task.

Rubric provenance differs by study, and the difference matters. The synthetic study derived its
rubric criteria from the same ground-truth fragments the composer indexes, which flatters any
condition carrying those fragments. The real studies use two independently authored families: a
\emph{practitioner} rubric written by a panel that saw only the work request, and a \emph{guidance}
rubric written by a panel that saw the request plus the oracle skills. Both are scored and both are
reported.

\subsection{Statistics}

Task-level scores are paired across conditions, so all contrasts are within-task differences. We
report the mean difference, a bias-corrected percentile bootstrap 95\% confidence interval over
20,000 resamples, the exact two-sided Wilcoxon signed-rank $p$ computed over the full null
distribution of sign assignments, Holm-corrected $p$ within each study's family of contrasts, the
paired standardised effect size $d_z$, the rank-biserial correlation, and the win/loss count over
the 12 tasks. Twelve tasks is a small sample and we resist reading unadjusted $p$ values.

Panel reliability is reported as exact agreement, quadratically weighted Cohen's $\kappa$ and
Krippendorff's $\alpha$ on the ordinal scale, all computed over the individual item scores. The
merge-precision comparison in Section~\ref{sec:s3} is a two-proportion $z$ test on independent
samples of 60 pairs each.

Token counts use a byte-pair encoder; a second encoder agreed to within a fraction of a percent on
sampled text. \emph{Injected context} is the guidance block alone. \emph{End to end} adds the task
framing and the model's own output.

\begin{table}[t]\centering\small
\begin{tabular}{llrrr}
\toprule
Study & Panels & Items & Weighted $\kappa$ & Krippendorff $\alpha$ \\
\midrule
1. Synthetic & 2 & 468 & 0.913 & 0.904 \\
2. Real library & 2 & 1{,}200 & 0.940 & 0.923 \\
3. Atoms & 2 & 960 & 0.934 & 0.916 \\
\bottomrule
\end{tabular}
\caption{Grader panel reliability. Exact agreement was 0.947, 0.955 and 0.951; every disagreement
in all three studies was within one point.}
\label{tab:rel}
\end{table}

\section{Strategy 1: fragment composition on a synthetic corpus}
\label{sec:s1}

\subsection{The method}

Skills are split at second-level headings into fragments; per-skill scaffolding sentences are
stripped by pattern. Fragments are embedded as L2-normalised TF-IDF vectors and clustered by
average-linkage agglomeration on cosine similarity. Each cluster is a competency, represented by
its shortest near-medoid member, which prefers the tersest phrasing of a shared obligation. For
each cluster a model is asked once, offline, for eight short phrasings of situations in which it
applies, plus six tags; the same is done for each skill description with fourteen phrasings. These
are indexed at double weight. The expansion model sees only the fragment text and was never shown
the evaluation tasks.

At request time, BM25 over skill descriptions and their expansions selects up to $k_{\max}$ skills
scoring at least $\alpha$ times the top score. Each competency $f$ is scored
\[
s(f) \;=\; w_r\,\widehat{\text{BM25}}(q, f) \;+\; w_o\,a(f) \;+\; w_s\,\frac{\log(1 + \sigma_f)}{\log(1 + \sigma_{\max})},
\]
where $a(f)$ is the normalised routing score of the best skill containing $f$ and $\sigma_f$ is the
number of skills it appears in. Competencies outside the routed set are kept at reduced weight, so
a routing miss degrades rather than deletes. Selection is greedy under a hard budget on the rendered
document.

\subsection{Decomposition quality}

The merging threshold was chosen on six skills and evaluated on the other six.

\begin{table}[h]\centering\small
\begin{tabular}{lrrrrr}
\toprule
Split & $k$ & true $k$ & ARI & Precision & Recall \\
\midrule
Development (6 skills) & 28 & 24 & 0.967 & 1.000 & 0.938 \\
Held out (6 skills) & 27 & 25 & 0.904 & 1.000 & 0.829 \\
Full corpus & 35 & 34 & 0.956 & 0.985 & 0.930 \\
\bottomrule
\end{tabular}
\caption{Cross-skill merging against ground truth at the development-selected threshold
$\tau = 0.14$. Precision is over pairwise co-clustering decisions, so 1.000 means no two distinct
competencies were ever merged.}
\label{tab:s1cluster}
\end{table}

Errors are one-sided: the method under-merges. That is the safer failure, since an unmerged
duplicate wastes budget while a false merge silently deletes an obligation.

\subsection{Coverage per token}

Before asking whether agents write better deliverables, a cheaper question is available: for a given
spend, how many of the competencies a task needs are present? On the eight tasks not used for
hyperparameter selection, the composite saturates at 98.2\% coverage for 1,272 tokens. Concatenating
the oracle skills reaches 98.4\% for 1,739 tokens, concatenating the routed skills 94.3\% for 2,489,
and loading the whole corpus 100\% for 10,137. At the 1,000-token operating point the composite
carries 0.87 covered competencies per thousand tokens against 0.57 for oracle concatenation and 0.38
for routed concatenation.

\subsection{Deliverable quality}

\begin{table}[t]\centering\small
\begin{tabular}{lrrrr}
\toprule
Condition & Rubric & Useful & Ctx & E2E \\
\midrule
A. No skill & 0.711 & 7.00 & 0 & 1{,}370 \\
D. Random fragments, token-matched & 0.820 & 7.58 & 1{,}074 & 2{,}496 \\
C. \sm{} composite & \textbf{0.948} & 8.46 & 1{,}098 & 2{,}623 \\
B. Whole skills, concatenated & 0.970 & 8.67 & 2{,}604 & 4{,}215 \\
\bottomrule
\end{tabular}
\caption{Strategy 1, synthetic corpus. Rubric is the fraction of the maximum; \emph{Useful} is the
1 to 10 usefulness judgement; \emph{Ctx} is injected context tokens and \emph{E2E} adds task
framing and model output. $n = 12$ tasks, two panels.}
\label{tab:s1main}
\end{table}

\begin{table}[t]\centering\small
\begin{tabular}{lrrrrrr}
\toprule
Contrast & $\Delta$ & 95\% CI & $d_z$ & $p$ & $p_{\text{Holm}}$ & W/L \\
\midrule
C $-$ A & $+0.237$ & $[+0.169, +0.303]$ & $+1.92$ & 0.0005 & 0.0024 & 12/0 \\
C $-$ D & $+0.128$ & $[+0.069, +0.194]$ & $+1.12$ & 0.0020 & 0.0059 & 10/0 \\
C $-$ B & $-0.022$ & $[-0.061, +0.018]$ & $-0.31$ & 0.3008 & 0.3008 & 2/7 \\
B $-$ A & $+0.259$ & $[+0.195, +0.322]$ & $+2.23$ & 0.0005 & 0.0024 & 12/0 \\
D $-$ A & $+0.109$ & $[+0.061, +0.159]$ & $+1.18$ & 0.0029 & 0.0059 & 10/1 \\
\bottomrule
\end{tabular}
\caption{Strategy 1 contrasts on the combined rubric. Exact Wilcoxon signed-rank, Holm-corrected
within the family of five.}
\label{tab:s1contrasts}
\end{table}

Three readings. The composite is not distinguishable from loading everything: the interval
$[-0.061, +0.018]$ excludes any deficit larger than 6.1 rubric points at 42\% of the injected
context. This is non-inferiority within that margin and not superiority, since B wins 7 of the 12
tasks by small amounts. Second, selection rather than brevity is what works: the random control
holds tokens constant and destroys only the mapping from task to content, and it loses 12.8 points
while winning no tasks. Third, the saving is real end to end, since responses are also shorter
(1,380 tokens against 1,467).

\subsection{Ablations}
\label{sec:s1abl}

At a 1,000-token budget on the held-out tasks, removing offline expansion costs 0.077 of coverage
and removing cross-skill merging costs 0.049; removing both costs 0.182, more than the sum, because
without merging the budget fills with duplicates and without expansion the wrong fragments are
ranked into it. The other two components fail. Removing the diversity rerank improves coverage by
0.025 and removing the support prior improves it by the same amount. Merging has already removed
the near-duplication the rerank exists to suppress, so its penalty only displaces relevant material;
the support prior is redundant with the routing term. Both are set to zero in later strategies.

\subsection{What Strategy 1 does not establish}

The corpus was built to contain the redundancy the method removes, and the rubric criteria were
derived from the same fragments the composer indexes. Both choices make the result measurable and
both flatter it. Strategy~2 repeats the experiment without either.

\section{Strategy 2: the same method on a production library}
\label{sec:s2}

\subsection{Reading real skills}

A real \texttt{SKILL.md} is prose interleaved with fenced code, tables of script invocations, ASCII
diagrams, argument placeholders such as \texttt{@\$1}, and links into \texttt{scripts/} and
\texttt{references/} that mean nothing outside the skill's own directory. Fragments are therefore
split into two classes at parse time. \emph{Guidance} fragments are prose obligations and are
mergeable. \emph{Pinned} fragments are anything whose value depends on the skill it came from: a
fenced block, a table, an invocation, or any text naming a file that ships with the skill. Pinned
fragments are never merged, are offered only when their own skill is routed, and are carried
verbatim with relative paths resolved to absolute ones.

Pinned material is 29.5\% of fragments and 35.7\% of corpus tokens. A composer that treats a skill
as a bag of sentences, which Strategy~1 did because its corpus contained nothing else, destroys a
third of a real library. From Strategy~2 onwards the composite is emitted as a valid
\texttt{SKILL.md}, so nothing downstream needs a bespoke format and the output can be re-indexed by
the reader that produced it.

\subsection{There is less redundancy than expected}

\begin{table}[h]\centering\small
\begin{tabular}{lrr}
\toprule
Merging policy & Clusters & Cross-skill \\
\midrule
None (raw fragments) & 344 & \\
Within-skill, raw text & 200 & 0 \\
Cross-skill, raw text & 186 & 24 \\
Cross-skill, normalised obligations & 220 & 57 \\
\bottomrule
\end{tabular}
\caption{Guidance fragments in the 45-skill library under each policy. Within-skill merging alone
removes 42\% of fragments.}
\label{tab:s2red}
\end{table}

Within-skill merging collapses 344 guidance fragments to 200, a factor of 1.72 that is real and
worth harvesting. Cross-skill merging on raw text adds 24 clusters across the entire library, and 6
inside the coherent ten-skill engineering sub-corpus. The synthetic corpus had a cross-skill
redundancy factor of 2.91 because we put it there. Real skills mostly repeat themselves.

\subsection{The normalisation trap}
\label{sec:s2trap}

Raw-text similarity misses semantic overlap. Two skills can both require a rehearsed reversal path
and share almost no vocabulary saying so. The available fix is to ask a model, offline and once, to
restate each fragment as a single domain-free obligation, then cluster those.

Against ground truth this works. Calibrated on six skills of the synthetic corpus and evaluated on
the other six, clustering normalised obligations reaches pairwise $F_1 = 0.953$ (ARI 0.952,
precision 0.911, recall 1.000) against 0.907 for raw text, whose recall is 0.829. Every true
co-clustering is recovered.

On the real library it fails silently. We sampled 120 cross-skill pairs, 60 merged by the
obligation clustering and 60 near-misses just below threshold, and gave them to eight independent
judges who were told nothing about how the pairs were chosen and asked whether a single well-written
instruction could replace both without losing anything a practitioner would act on differently.
Four of 60 merges were endorsed, none of the 42 high-confidence judgements among them, and mean
judge confidence was 4.43 of 5.

Inspecting the rejects explains it. Normalisation was instructed to drop domain nouns, tool names
and examples and keep the imperative. On a heterogeneous library that turns ``use the design
philosophy as the foundation for a single highly visual page'' and ``connector results arrive inline,
so search wide and keep deliberately'' into two generic sentences about being deliberate, which then
cluster. The information that distinguished them is the information the normalisation removed.
Fitting a conjunctive rule on obligation similarity and raw similarity together, selected on half
the panel's labels and reported on the other half, reaches precision 0.111, because with four
genuine duplicates in the sampled band no threshold can do better. Cross-skill merging is therefore
off by default from Strategy~2 onwards, enabled only when a caller asserts a coherent single-domain
corpus.

\subsection{Routing is the ceiling}

Skill routing over 45 skills reaches $F_1 = 0.674$ (precision 0.681, recall 0.667) against the
oracle sets, close to the 0.662 the same router reached over 12. Three attempted improvements are
worth recording because none worked. Aggregating each skill's best fragment scores into its routing
score, a standard passage-evidence move, was neutral at every weight. Setting BM25's length
normalisation $b$ toward zero, on the argument that description length in this format tracks scope
rather than verbosity, fixed individual cases without moving aggregate $F_1$. Hedging the routing
width, narrow when the score distribution is peaked and wide when it is flat, raised recall from
0.667 to 0.708 at a precision cost, and we kept it because fragment selection filters afterwards.

The residual failures are not tunable. One task describes a nightly job that ``produced silently
wrong totals \dots{} runs fine, exits zero''. The \texttt{debug} skill carries the offline expansion
``output is wrong and i dont know why'', which is semantically exact and lexically disjoint. No
weighting of term statistics bridges that, and an embedding router is where the next gain lies.

What we could add is a way to fail quietly. The top skill's raw combined score, calibrated only on
in-corpus probes (each skill's own expansions, short and concatenated-long) against out-of-corpus
probes about cooking, football and gardening, gives a floor at which the server returns nothing. It
keeps 95\% of in-corpus queries and abstains on 14 of 15 out-of-corpus ones. All 12 evaluation tasks
clear it, and the two lowest are the two the router gets wrong.

\subsection{Deliverable quality}

\begin{table}[t]\centering\small
\begin{tabular}{lrrrrr}
\toprule
Condition & Pract. & Guid. & Both & Ctx & Recov. \\
\midrule
A. No skill & 0.910 & 0.694 & 0.802 & 0 & \\
D. Random, token-matched & 0.894 & 0.658 & 0.776 & 873 & $-23\%$ \\
C. \sm{} composite & 0.871 & 0.808 & 0.840 & 967 & 33\% \\
E. \sm{}, token-matched to B & 0.900 & 0.810 & 0.855 & 2{,}078 & 47\% \\
B. Whole skills, oracle-selected & 0.873 & 0.956 & \textbf{0.915} & 2{,}414 & 100\% \\
\bottomrule
\end{tabular}
\caption{Strategy 2, real library. \emph{Pract.} is the practitioner rubric, \emph{Guid.} the
guidance rubric. \emph{Recov.} is the share of B's improvement over A that the condition recovers.
Condition B is given oracle skill selection, which the composite is not, so the comparison is
conservative.}
\label{tab:s2main}
\end{table}

\begin{table}[t]\centering\small
\begin{tabular}{lrrrrrr}
\toprule
Contrast & $\Delta$ & 95\% CI & $d_z$ & $p$ & $p_{\text{Holm}}$ & W/L \\
\midrule
C $-$ B & $-0.075$ & $[-0.113, -0.033]$ & $-1.00$ & 0.0059 & 0.0293 & 1/11 \\
E $-$ B & $-0.059$ & $[-0.106, -0.014]$ & $-0.70$ & 0.0293 & 0.1172 & 1/9 \\
E $-$ C & $+0.016$ & $[-0.021, +0.051]$ & $+0.24$ & 0.4775 & 0.7607 & 6/5 \\
C $-$ A & $+0.038$ & $[-0.029, +0.109]$ & $+0.29$ & 0.3804 & 0.7607 & 8/4 \\
C $-$ D & $+0.064$ & $[-0.002, +0.132]$ & $+0.51$ & 0.0962 & 0.2886 & 9/3 \\
B $-$ A & $+0.113$ & $[+0.056, +0.177]$ & $+1.00$ & 0.0020 & 0.0117 & 10/1 \\
\bottomrule
\end{tabular}
\caption{Strategy 2 contrasts, combined rubric. Only the two contrasts against whole skills survive
Holm correction.}
\label{tab:s2contrasts}
\end{table}

The Strategy~1 headline does not replicate. The composite gives up 7.5 rubric points and loses 11
of 12 tasks. Condition E raises the ceiling to the baseline's own token cost and lifts the relevance
floor; it gains $+0.016$, indistinguishable from zero, and still loses to whole skills. The deficit
is in what gets selected, not in how much room it has.

Where the two rubric families disagree is the informative part. On the practitioner rubric, written
by a panel that saw only the request, \emph{no skill} scores highest at 0.910: a capable model
already writes a competent pull-request review. Skills earn their place on the guidance rubric,
where whole skills reach 0.956 against 0.694 for no skill, and the composite reaches 0.808, about
half the distance. A skill's value sits in its specifics, and specifics are what fragment selection
drops first. Strategy~1, whose rubrics were all of the guidance kind, saw only the flattering half
of this picture.

The genre split is usable. On document production, where skills carry script paths, library
gotchas and format-specific procedure the model does not otherwise have, the composite recovers 51\%
of the benefit at 37\% of the cost and beats both no guidance (0.833 against 0.733) and the random
control (0.688). On engineering tasks it recovers nothing: 0.845 against 0.852 for no skill at all,
because those skills largely restate practice the model already has and average 632 tokens, leaving
nothing to save. Composition pays where a skill teaches and not where it reminds.

The random control is also informative. It scores below no guidance at all, recovering $-23\%$ of
the benefit, so irrelevant procedural text is worse than silence.

\section{Strategy 3: obligation and lesson atoms}
\label{sec:s3}

\subsection{The representation}

Strategy~2's merging failure had a precise shape. To compare guidance across skills it normalised
each fragment to a domain-free obligation, and to keep the budget it then kept one fragment per
cluster and dropped the rest. The first step made unlike things look alike; the second made that
permanent.

Strategy~3 stops conflating two things a fragment contains. An \textbf{obligation} is one imperative
sentence: what the practitioner must do and when, in plain verbs with tool, product and file names
stripped. It is the only field ever compared between skills. A \textbf{lesson} is the substance: the
reason, the mechanism, the exact command, the failure it prevents, the names and numbers and paths,
kept in the source's own words and verbatim where the source was verbatim. A lesson is never merged,
never paraphrased, and never dropped when its obligation merges. A merged atom carries one obligation
and every contributing lesson, each labelled with its source skill.

The question ``can one instruction replace both?'', which Strategy~2 had to answer yes to and could
not, becomes ``is this obligation a truthful heading for both lessons?'', which is weaker and far more
often true.

Skills cannot be used in this form as written, so Strategy~3 adds a reformatter: an offline pass in
which a model reads one \texttt{SKILL.md} and emits its atoms, reproducing code blocks, command
tables and paths exactly inside the lessons they belong to. Refactoring the 45 skills produced 1,134
atoms, 393 of them pinned, 20,311 tokens of obligation and 90,726 of lesson. The atomised library is
1.29 times the size of the source, an index-time cost paid once; obligations are 18\% of it.

\subsection{What refactoring costs}

\begin{table}[h]\centering\small
\begin{tabular}{lr}
\toprule
Measure & Value \\
\midrule
Instructions preserved & 98.8\% of 588 \\
Code blocks and command tables preserved & 100\% \\
File paths preserved & 94.8\% \\
Distortions (atoms state what the source does not) & 49 \\
Skills showing at least one distortion & 9 of 11 \\
\bottomrule
\end{tabular}
\caption{Reformatting fidelity, audited by independent agents reading the original and the atoms
side by side.}
\label{tab:s3fid}
\end{table}

Substance survives. Modality does not. Nearly every distortion was of one kind: the source's hedge
became the atom's requirement. ``Prefer using the imperative form'' became ``Write every instruction in
the imperative form''. A skill that says to ask when no mode is given acquired an atom demanding a
mode. This is a predictable consequence of asking for imperatives, since obligations are
grammatically mandatory and prose is not, so a rewrite into obligations is a rewrite into
mandatoriness. A reformatter should be constrained to preserve must, should and may, and forbidden
from introducing a requirement the source does not impose, and it should still be audited, because
the pressure is structural.

One class of distortion is removable without a model. The reformatter emitted a bare
\texttt{recalc.py} where the skill ships \texttt{scripts/recalc.py}. An invented path reads as
authoritative and fails only when someone runs it, so every reference is now checked against files
that exist, resolved when unambiguous and dropped otherwise: 135 kept, 110 dropped. Many of the
drops are patterns rather than paths, such as \texttt{ppt/slides/slideN.xml}, so this is a
conservative filter rather than an error count.

\subsection{Merging, re-adjudicated}

Under the atom design nothing is replaced by merging, so the corrected question asks whether the
shared obligation is a truthful heading for both lessons, telling the judge explicitly that the two
lessons need not resemble each other and that both are kept. An independent panel of eight judges
saw 120 pairs on the same sampling design as Strategy~2.

\begin{table}[h]\centering\small
\begin{tabular}{lrr}
\toprule
& Strategy 2 & Strategy 3 \\
\midrule
Merged pairs endorsed & 4/60 & \textbf{49/60} \\
Precision & 0.067 & \textbf{0.817} \\
Precision, judge confidence $\geq 4$ & 0.000 & 0.808 \\
Near-misses the panel would merge & 0/60 & 5/60 \\
Band-limited recall & & 0.907 \\
\bottomrule
\end{tabular}
\caption{Cross-skill merge quality on the same 45-skill library, independent panels, mean judge
confidence 4.5 of 5. Difference in precision $+0.750$, 95\% CI $[+0.634, +0.866]$, $z = 8.3$,
$p < 0.001$.}
\label{tab:s3adj}
\end{table}

The merges are the ones a practitioner would want. \emph{Obtain the person's explicit approval
before the final irreversible action} is stated once and carries the lessons of eight consumer-task
skills. \emph{Assume the standard packages are installed and install one only after an import fails}
carries the specific package, command and footgun list from \texttt{docx}, \texttt{pptx} and
\texttt{xlsx}.

\subsection{Deliverable quality, and the ablation that explains it}

\begin{table}[t]\centering\small
\begin{tabular}{lrrrr}
\toprule
Condition & Pract. & Guid. & Both & Ctx \\
\midrule
B. Whole skills, oracle-selected & 0.890 & 0.950 & \textbf{0.920} & 2{,}414 \\
C. Strategy 2 fragment composite & 0.867 & 0.808 & 0.838 & 967 \\
F. Strategy 3, obligation and lesson & 0.844 & 0.750 & 0.797 & 1{,}230 \\
G. Strategy 3, lessons only & 0.873 & 0.762 & 0.818 & 1{,}309 \\
\bottomrule
\end{tabular}
\caption{Strategy 3. Condition F is worse than the fragment composite it was designed to replace,
despite resting on merges twelve times more likely to be correct.}
\label{tab:s3main}
\end{table}

\begin{table}[t]\centering\small
\begin{tabular}{lrrrrrr}
\toprule
Contrast & $\Delta$ & 95\% CI & $d_z$ & $p$ & $p_{\text{Holm}}$ & W/L \\
\midrule
G $-$ F & $+0.021$ & $[-0.007, +0.050]$ & $+0.40$ & 0.2500 & 0.6094 & 6/5 \\
G $-$ C & $-0.020$ & $[-0.071, +0.031]$ & $-0.21$ & 0.5049 & 0.6094 & 5/6 \\
F $-$ C & $-0.041$ & $[-0.095, +0.007]$ & $-0.43$ & 0.2031 & 0.6094 & 3/7 \\
G $-$ B & $-0.102$ & $[-0.147, -0.060]$ & $-1.31$ & 0.0024 & 0.0098 & 1/11 \\
C $-$ B & $-0.082$ & $[-0.106, -0.058]$ & $-1.84$ & 0.0005 & 0.0024 & 0/12 \\
\bottomrule
\end{tabular}
\caption{Strategy 3 contrasts, combined rubric. After Holm correction only the contrasts against
whole skills are distinguishable; the three composites are not separable from one another.}
\label{tab:s3contrasts}
\end{table}

Solving the representation problem did not move the outcome. Condition F loses 4.1 rubric points to
the fragment composite and a full point of rated usefulness on the 1 to 10 scale
($-1.000$, exact Wilcoxon $p = 0.008$ unadjusted), though on the combined rubric that difference
does not survive correction.

If the atoms are right and the composite is worse, the suspect is what the composite prints, so we
rendered the identical selection with the obligation headings removed and the freed budget spent on
further lessons. Condition G recovers $+0.021$ on the rubric and $+0.58$ on rated usefulness (7.33
against 6.75), landing statistically level with the fragment composite. Printing the obligations
was costing quality.

The explanation is Strategy~2's genre result turned on this design. Composition pays where a skill
teaches and not where it reminds, because a capable model already holds the reminders. An obligation
is by construction the reminder half: it is what remains once the specifics are stripped. The atom
design spends 18\% of the atomised library, and a comparable share of every composite's budget,
restating in general terms things the model was going to do anyway, in the hardened modality of
Table~\ref{tab:s3fid}, which displaces specifics with commands.

We note the size of the effect honestly. $+0.021$ with an interval that includes zero is weak
evidence on its own. It is supported by the direction being consistent across two measures, by the
mechanism being independently motivated, and by G landing level with C while F does not.

\section{The retained system}
\label{sec:retained}

We retain Strategy~3's representation and reject its rendering. The shipped configuration indexes
the library as obligation and lesson atoms, merges across skills on obligations only, keeps every
contributing lesson, and emits lessons alone, holding the obligation in the selection trace where a
person debugging a composition can see it. Three guards from Strategy~2 stay, each of them the
result of something that went wrong.

\textbf{Abstention.} Out-of-corpus queries previously produced confident irrelevance. The server now
returns nothing when no skill scores above the calibrated floor, which catches 14 of 15 negatives
while keeping 95\% of positives.

\textbf{Never worse than the source.} A composite is capped at 90\% of the cost of the skills it
drew on, so composing cannot cost more than loading them. Without it, on small skills, the greedy
selector filled a generous budget with unrelated material.

\textbf{Pinned material carried verbatim.} Code, tables, invocations and file pointers are never
merged or paraphrased, and relative paths are resolved and validated against files that exist.

The composite is emitted as a valid \texttt{SKILL.md}, so no consumer needs a bespoke format and the
output can be re-indexed by the reader that produced it. Composition on the request path stays
deterministic and model-free: it is a pure function of the index, the query and the configuration,
which makes it reproducible, auditable back to its source skills, and cheap enough to call before
every task.

We recommend this configuration on two grounds, and not on a third. It reduces injected context by
46\% to 60\% at a quality cost that is real and bounded, and it makes a library de-duplicable and
searchable: one obligation, every skill that imposes it, every lesson attached. It does not improve
deliverables, and we found no configuration that did.

\section{Availability}
\label{sec:avail}

The implementation, the twelve-skill synthetic corpus, every results file and the scripts that
regenerate the tables in this article are at \url{https://github.com/mathieubrglt/skillmerge} under
the MIT License. This article is available under CC BY 4.0. Running
\texttt{tools/verify\_article.py} checks all 194 reported numbers against the stored results and
exits non-zero on any mismatch; \texttt{tools/stats\_all.py} regenerates every contrast, effect size
and reliability coefficient.

The 45-skill production library is not redistributed. Seventeen of those skills carry their own
licence header and six are marked proprietary, and the atoms are close derivatives of them by
design, since a lesson reproduces its source verbatim. The repository instead carries
\texttt{docs/SKILLS-MANIFEST.json}, which records every skill the study read with the SHA-256 of the
document as it stood, so a reader with the same installation can confirm they are looking at the
same corpus before rebuilding the index locally. The synthetic corpus is regenerated byte-identically
from its seeded generator, so the calibration results in Sections~\ref{sec:s1} and \ref{sec:s2trap}
can be reproduced with nothing else installed.

\section{Threats to validity}

\textbf{Sample size.} Twelve tasks per study. Confidence intervals are over tasks, not over corpora
or domains, and after Holm correction the differences among composites in Strategy~3 are not
distinguishable. We rest the retained recommendation on the merge-precision result, which is a
two-proportion comparison on 120 independently judged pairs, rather than on the composite contrasts.

\textbf{Rubric and judge provenance.} Both rubric families, both grader panels and both adjudication
panels are language models. Reliability between panels is high (Table~\ref{tab:rel}) and the
adjudication panels were independent of the systems they judged and blind to how pairs were
selected, but agreement is not correctness, and graders may share blind spots with the systems they
grade.

\textbf{The adjudication sample is band-limited.} It covers pairs the clustering surfaced or nearly
surfaced, so the recall figure of 0.907 is recall within reach of the lexical blocking rather than
recall against all true duplicates in the library.

\textbf{One reformatter, one prompt.} The modality hardening in Table~\ref{tab:s3fid} is a property
of that prompt as much as of the idea, and a reformatter constrained to preserve modality might
produce atoms that cost less quality. We did not test one.

\textbf{Composition destroys order.} A skill's sequence carries meaning and the four-phase regrouping
is a crude replacement. On procedural skills this should matter more than on checklist-shaped ones,
and we do not separate the two.

\textbf{Lexical retrieval throughout.} No embedding model was available in the environment, so
routing is BM25 with offline query expansion. Routing recall of 0.71 caps every result here, and the
one failure we can characterise precisely is a vocabulary gap that no term weighting closes.

\section{Conclusion}

The composition idea is sound on a corpus built to contain the redundancy it removes and does not
survive contact with a library written by people who were not thinking about composition. What
survives is narrower and still useful: obligations are a good index key and a poor payload. Split
each unit into the imperative and the teaching, compare only the imperative, keep every piece of
teaching, and a heterogeneous skill library becomes de-duplicable at 0.817 precision where the
previous method managed 0.067. Then print only the teaching, because the model already has the
imperatives.

The practical consequence for anyone maintaining a skill library is that the obligation index is
probably worth more at authoring time than at request time. It shows exactly where skills overlap,
which is the question a maintainer has and cannot currently answer. The runtime case for composition
rests on context cost and on abstention, and anyone promising quality parity with loading the skills
outright should be asked for their rubric.

\bibliographystyle{plainnat}
\begin{thebibliography}{10}
\providecommand{\natexlab}[1]{#1}

\bibitem[Anthropic(2024)]{anthropic2024mcp}
Anthropic.
\newblock Introducing the {M}odel {C}ontext {P}rotocol, 2024.
\newblock \url{https://www.anthropic.com/news/model-context-protocol}.

\bibitem[Anthropic(2025)]{anthropic2025skills}
Anthropic.
\newblock Agent {S}kills, 2025.
\newblock \url{https://docs.claude.com/en/docs/agents-and-tools/agent-skills}.

\bibitem[Carbonell and Goldstein(1998)]{carbonell1998mmr}
J.~Carbonell and J.~Goldstein.
\newblock The use of {MMR}, diversity-based reranking for reordering documents and producing
  summaries.
\newblock In \emph{SIGIR}, pages 335--336, 1998.

\bibitem[Jiang et~al.(2023)]{jiang2023llmlingua}
H.~Jiang, Q.~Wu, C.-Y. Lin, Y.~Yang, and L.~Qiu.
\newblock {LLMLingua}: Compressing prompts for accelerated inference of large language models.
\newblock In \emph{EMNLP}, 2023.
\newblock \url{https://arxiv.org/abs/2310.05736}.

\bibitem[Lewis et~al.(2020)]{lewis2020rag}
P.~Lewis, E.~Perez, A.~Piktus, F.~Petroni, V.~Karpukhin, N.~Goyal, H.~K\"uttler, M.~Lewis,
  W.~Yih, T.~Rockt\"aschel, S.~Riedel, and D.~Kiela.
\newblock Retrieval-augmented generation for knowledge-intensive {NLP} tasks.
\newblock In \emph{NeurIPS}, 2020.
\newblock \url{https://arxiv.org/abs/2005.11401}.

\bibitem[Liu et~al.(2024)]{liu2024lost}
N.~F. Liu, K.~Lin, J.~Hewitt, A.~Paranjape, M.~Bevilacqua, F.~Petroni, and P.~Liang.
\newblock Lost in the middle: How language models use long contexts.
\newblock \emph{Transactions of the ACL}, 12:157--173, 2024.
\newblock \url{https://arxiv.org/abs/2307.03172}.

\bibitem[Nogueira et~al.(2019)]{nogueira2019doc2query}
R.~Nogueira, W.~Yang, J.~Lin, and K.~Cho.
\newblock Document expansion by query prediction, 2019.
\newblock \url{https://arxiv.org/abs/1904.08375}.

\bibitem[Pan et~al.(2024)]{pan2024llmlingua2}
Z.~Pan, Q.~Wu, H.~Jiang, M.~Xia, X.~Luo, J.~Zhang, Q.~Lin, V.~R\"uhle, Y.~Yang, C.-Y. Lin,
  H.~V. Zhao, L.~Qiu, and D.~Zhang.
\newblock {LLMLingua-2}: Data distillation for efficient and faithful task-agnostic prompt
  compression, 2024.
\newblock \url{https://arxiv.org/abs/2403.12968}.

\bibitem[Robertson and Zaragoza(2009)]{robertson2009bm25}
S.~Robertson and H.~Zaragoza.
\newblock The probabilistic relevance framework: {BM25} and beyond.
\newblock \emph{Foundations and Trends in Information Retrieval}, 3\penalty0 (4):\penalty0
  333--389, 2009.

\end{thebibliography}

\appendix
\section{Reproducibility}

Every artefact is deterministic given the seed. \texttt{tools/build\_corpus.py} regenerates the
synthetic corpus and its ground truth from the fragment library. \texttt{tools/eval\_cluster\_split.py}
reproduces Table~\ref{tab:s1cluster}, \texttt{tools/sweep\_budget.py} the coverage figures,
\texttt{tools/ablations.py} Section~\ref{sec:s1abl}, \texttt{tools/calibrate\_v2.py} the calibration
in Section~\ref{sec:s2trap}, \texttt{tools/analyze\_merges.py} and
\texttt{tools/analyze\_v3\_merges.py} the two adjudications, \texttt{tools/analyze\_fidelity.py}
Table~\ref{tab:s3fid}, and \texttt{tools/stats\_all.py} every contrast and reliability coefficient in
this article. \texttt{tools/verify\_paper.py}, \texttt{verify\_paper2.py} and
\texttt{verify\_paper3.py} check the reported numbers against the stored results and fail on
mismatch.

Selection protocol. The synthetic study's merging threshold was chosen on six skills and reported on
the other six; its composition hyperparameters were chosen on four development tasks and its
coverage and ablation figures reported on the remaining eight. The abstention floor was calibrated
only on in-corpus and out-of-corpus probes, never on an evaluation task. The conjunctive merge rule
of Section~\ref{sec:s2trap} was fitted on half the adjudication labels and reported on the other
half. Offline expansions and atoms were generated by agents given only skill text, with instructions
not to read anything else in the repository; they had no access to the task suites.

\end{document}
